\section{Conclusion}
\label{sec:conclusion}

We presented \aris{}, a pipeline for automatically constructing training-ready
synthetic datasets for image editing tasks.
The central contribution is the free-form VLM-guided rendering evolution loop,
in which an AI agent reads natural language feedback from a vision-language model reviewer
and selects rendering parameter adjustments from a structured action space,
iterating until the reviewer confirms acceptable quality.
This approach replaces manual Blender parameter tuning with an automated feedback cycle,
enabling scalable and consistent 3D dataset construction.

We instantiated \aris{} for rotation-conditioned image editing,
producing ARIS-Rotate: a 50-object, 350-pair dataset
with natural-language viewpoint instructions.
Ablation experiments confirmed that the evolution loop improves rendering quality
over single-pass rendering.
LoRA fine-tuning of Qwen Image Edit 2511 on ARIS-Rotate
demonstrated measurable improvement over the base model on the test set,
validating the downstream utility of automatically constructed data.

\paragraph{Future work.}
Several directions extend this work naturally.
First, expanding the action space to cover failure modes currently unresolved
(global tone mapping, background replacement) would improve coverage.
Second, applying \aris{} to additional editing tasks---lighting variation,
material substitution, scale change---would validate generality.
Third, scaling to larger object sets (500+) with broader category coverage
would provide the data volume needed for training stronger editing models.
Finally, integrating an adaptive round budget that allocates more evolution rounds
to harder objects would improve efficiency on large-scale runs.
